% Draft subsection extending "Anchored Langevin Algorithms" (arXiv:2509.19455).
% Numbering placeholders (\ref{...}) point at the existing labels of that paper:
%   eq:al-sde  (Eq. 5-7)   thm:invariant (Thm 2)   thm:heavy (Thm 3)
%   thm:timechange (Thm 11)   asm:msq (Asm 12)   thm:w2 (Thm 14)
% All numbers below are produced by the code at
%   https://github.com/pervezali1/Accelerating-Non-Differentiable-sampling

\section{Non-reversible anchored Langevin dynamics}
\label{sec:skew}

The anchored Langevin SDE \eqref{eq:al-sde} is reversible with respect to
$\pi\propto e^{-U}$.  Reversibility is not required for correctness, and it is
known that a suitable non-reversible perturbation can only improve the rate of
convergence of a diffusion to its stationary law
\citep{hwang1993accelerating,hwang2005accelerating,lelievre2013optimal,%
duncan2016variance,reybellet2015irreversible}.  We show that such a
perturbation is available inside the anchored framework, that it costs one
matrix--vector product per step, and that it is effective precisely on the
anisotropic heavy-tailed targets that the reversible dynamics handles worst.

Let $J\in\mathbb R^{d\times d}$ be a constant skew-symmetric matrix,
$J^\top=-J$, and consider
\begin{equation}\label{eq:sal-sde}
  dX_t \;=\; e^{(U-U_0)(X_t)}\,(J-I)\,\nabla U_0(X_t)\,dt
        \;+\;\sqrt2\,e^{(U-U_0)(X_t)/2}\,dW_t ,
\end{equation}
which we call the \emph{skew-anchored Langevin SDE}.  Setting $J=0$ recovers
\eqref{eq:al-sde}.

\subsection{Invariance and ergodicity}

\begin{theorem}\label{thm:skew-invariant}
Let $U_0\in C^2$, $U\in C^1$ with $\int e^{-U}<\infty$, and suppose
\eqref{eq:sal-sde} does not explode.  Then $\pi\propto e^{-U}$ is an invariant
measure of \eqref{eq:sal-sde} for every constant skew-symmetric $J$.
\end{theorem}

\begin{proof}[Proof sketch]
The generator is
$\mathcal L_Jf=e^{U-U_0}\big(\Delta f-\langle\nabla U_0,\nabla f\rangle
+\langle J\nabla U_0,\nabla f\rangle\big)$.  Using
$e^{(U-U_0)}e^{-U}=e^{-U_0}$, the added drift $c=e^{U-U_0}J\nabla U_0$ satisfies
\[
  \nabla\!\cdot\!\big(e^{-U}c\big)=\nabla\!\cdot\!\big(e^{-U_0}J\nabla U_0\big)
  =e^{-U_0}\Big(\operatorname{Tr}\big(J\nabla^2U_0\big)-\langle\nabla U_0,J\nabla U_0\rangle\Big)=0,
\]
since $\operatorname{Tr}(JH)=0$ for skew $J$ and symmetric $H$, and
$\langle v,Jv\rangle=0$.  Hence $\int\mathcal L_Jf\,d\pi=0$ on $C_c^\infty$.
\end{proof}

In $L^2(\pi)$ the added term is antisymmetric, so the Dirichlet form of
Lemma~\ref{lem:dirichlet} is unchanged; the perturbation cannot slow
convergence and generically accelerates it.

\begin{proposition}\label{prop:skew-lyapunov}
Let $V=\Psi\circ U_0$ with $\Psi\in C^2$ increasing.  Then
$\mathcal L_JV=\mathcal LV$ identically.  Consequently every drift condition of
the form of Assumption~\ref{asm:drift} verified with such a $V$ holds for every
skew $J$ with the same constants, and Theorems~\ref{thm:invariant} and
\ref{thm:heavy} extend verbatim.
\end{proposition}

\begin{proof}
$\langle J\nabla U_0,\nabla V\rangle=\Psi'(U_0)\langle J\nabla U_0,\nabla U_0\rangle=0$.
\end{proof}

For the heavy-tailed targets of Theorem~\ref{thm:heavy}, where
$U_0=\beta\log q$, Proposition~\ref{prop:skew-lyapunov} applies with
$\Psi(u)=e^{u/\beta}$, so $V=q$.  The random time change of
Section~\ref{sec:timechange} also carries over: $X_t=Z_{\ell(t)}$ where $Z$
solves $dZ_t=(J-I)\nabla U_0(Z_t)dt+\sqrt2\,d\widetilde W_t$ and $\ell$ is
unchanged, and the two Euler--Maruyama schemes remain pathwise identical under
synchronous coupling.

\begin{remark}[The mean-square analysis does not extend]
Assumption~\ref{asm:msq} asks for a one-sided Lipschitz constant $m>\alpha$.
For the Student-$t$ target of Example~\ref{ex:studentt} one finds
$m_J=\tfrac{2\beta}{\nu}\lambda_{\min}(\Sigma^{-1}-S)$ with
$S=\tfrac12[J,\Sigma^{-1}]$, so $m_J>0$ requires
$\rho_J:=\lambda_{\max}(\Sigma^{1/2}S\Sigma^{1/2})<1$; and $\rho_J=0$ exactly
when $J$ commutes with $\Sigma$, which is also exactly when $J$ cannot
accelerate.  Moreover $\alpha<m$ already fails at $J=0$ unless
$\beta>\tfrac d2\kappa(\Sigma)$, the condition of
Corollary~\ref{cor:studentt}.  We therefore do not extend Theorem~\ref{thm:w2};
for the canonical anchor $\beta=\tfrac{d+\nu}{2}-1$ we instead compute the
second moment of the discretisation exactly (Section~\ref{sec:skew-msq}).
\end{remark}

\subsection{Exact second-moment analysis of the discretisation}
\label{sec:skew-msq}

For $U=\tfrac{d+\nu}{2}\log q$ and $U_0=\beta\log q$ with
$\beta=\tfrac{d+\nu}{2}-1$, the drift of \eqref{eq:sal-sde} is linear and the
diffusion coefficient is the scalar $q^{1/2}$, so the Euler--Maruyama scheme is
\[
  x_{k+1}=(I-\eta B)x_k+\sqrt{2\eta}\,q(x_k)^{1/2}\xi_{k+1},
  \qquad B=\tfrac{2\beta}{\nu}(I-J)\Sigma^{-1}.
\]
Because $q$ is exactly quadratic and the noise is isotropic given $x$, the
second-moment matrix obeys the \emph{exact} affine recursion
\begin{equation}\label{eq:msq-recursion}
  C_{k+1}=MC_kM^\top+2\eta\Big(1+\tfrac{\operatorname{Tr}(\Sigma^{-1}C_k)}{\nu}\Big)I,
  \qquad M=I-\eta B,
\end{equation}
whose linear part is
$L=M\otimes M+\tfrac{2\eta}{\nu}\operatorname{vec}(I)\operatorname{vec}(\Sigma^{-1})^\top$.
The scheme is mean-square stable iff $\rho(L)<1$; the stationary covariance of
the chain is the fixed point of \eqref{eq:msq-recursion}, giving the
discretisation bias exactly rather than as an $O(\eta)$ bound; and
$-\log\rho(L)$ is the exact per-iteration rate.  Note that the stationary
covariance of the \emph{SDE} is $\tfrac{\nu}{\nu-2}\Sigma$ for every $J$, since
the $J$ terms cancel in $BC+CB^\top$; this is what makes a comparison across
$J$ at equal bias well posed.

\subsection{Choosing $J$, and a ceiling}

Since $\operatorname{Tr}(JA)=0$ for skew $J$ and symmetric $A$, the eigenvalues
of $(I-J)A$ sum to $\operatorname{Tr}(A)$, whence
\begin{equation}\label{eq:ceiling}
  \min_i\operatorname{Re}\lambda_i\big((I-J)A\big)\le\frac{\operatorname{Tr}(A)}{d},
\end{equation}
against $\lambda_{\min}(A)$ at $J=0$.  The bound \eqref{eq:ceiling} is attained:
$(I+K)A$ is similar to $A+A^{1/2}KA^{1/2}$, so it suffices to rotate to a basis
in which $A$ has constant diagonal $\bar a=\operatorname{Tr}(A)/d$ and there
replace the strictly lower triangle by minus its transpose, a skew move leaving
an upper-triangular matrix with spectrum $\{\bar a\}$.

\begin{proposition}[No gain on radial targets]\label{prop:radial}
If $U$ and $U_0$ are radial, then for every constant skew $J$ the added drift
generates rotations, which commute with $\mathcal L$ and act unitarily on
$L^2(\pi)$.  Hence $\|e^{t\mathcal L_J}f\|_{L^2(\pi)}=\|e^{t\mathcal L}f\|_{L^2(\pi)}$
for all $f$ and $t$: the perturbation has no effect.
\end{proposition}

Proposition~\ref{prop:radial} applies to the target
$\pi\propto(1+\|x\|^2)^{-\iota}$ of Section~\ref{sec:heavyexp}; note also that
$\iota>1+d/2$ with $\iota=2$ forces $d=1$, where the only skew matrix is $0$.
The experiments below therefore use anisotropic heavy-tailed targets.

\subsection{Numerical experiments}

We take $\pi$ to be the $d$-dimensional Student-$t$ with $\nu$ degrees of
freedom and scale $\Sigma$ of condition number $\kappa$, anchored with
$\beta=\tfrac{d+\nu}{2}-1$, and follow the protocol of
Section~\ref{sec:heavyexp}: $5{,}000$ particles, priors $\mathcal N(0,10I)$ and
$\mathrm{Uniform}(-5,5)$, sliced 2-Wasserstein distance against the exact
projected quantiles.  Since a larger $\|J\|$ forces a smaller stable stepsize,
each variant is run at the stepsize that puts its stationary covariance at a
common relative error of $2\%$, computed from \eqref{eq:msq-recursion}; the
comparison is then at equal accuracy and equal cost per iteration.

Table~\ref{tab:skew-speedup} reports the resulting speed-up
$-\log\rho(L_J)\,/\,{-\log\rho(L_0)}$.

\begin{table}[t]
\centering
\caption{Per-iteration speed-up at equal discretisation bias, exact.}
\label{tab:skew-speedup}
\begin{tabular}{lrrrr}
\toprule
$\kappa(\Sigma)$ & $d=2$ & $d=3$ & $d=5$ & $d=10$\\
\midrule
1    & 1.00 & 1.00 & 1.00 & 1.00\\
10   & 1.20 & 1.19 & 1.13 & 1.16\\
100  & 6.83 & 3.57 & 2.08 & 1.53\\
1000 & 64.2 & 20.3 & 9.49 & 4.91\\
\bottomrule
\end{tabular}
\end{table}

Simulation confirms the prediction.  For $d=2$, $\nu=5$, $\kappa=100$ averaged
over $10$ runs, the iterations needed to bring the variance along the slowest
direction of $\Sigma$ within $10\%$ of the truth fall from $3316$ at $J=0$ to
$831$ at $\|J\|_2=3.46$, a factor $4.0$; no run diverged.

Two caveats deserve emphasis.  First, the gain in the \emph{continuous-time}
rate is much larger than the realisable one --- up to $380\times$ against
$64\times$ per iteration --- because the larger drift shrinks the mean-square
stable stepsize; optimising \eqref{eq:ceiling} alone therefore selects a
$\|J\|$ far beyond the useful range.  Second, the gain decreases with dimension
at fixed $\kappa$, and we make no claim about the high-dimensional regime.

Finally, we combine both difficulties the anchored framework is designed for by
taking $U=\iota\log q+\sum_i p_\lambda(x_i)$ with $p_\lambda$ the MCP penalty of
Section~\ref{sec:regularisers} and $U_0=\beta\log q+\sum_i p^\varepsilon_\lambda(x_i)$
its smoothing \eqref{eq:mcp-smoothed}.  The penalty is bounded, so the
polynomial tail is preserved and exact reference samples are available by
rejection from the Student-$t$ core.
